\documentclass[11pt, a4paper]{article}

% Lecture notes with numbered definitions, theorems and examples, and a
% boxed key point. Change the course and lecture details in the title, then
% replace each section with your own material.

% Fonts and encoding: pdfLaTeX needs these; XeLaTeX and Tectonic do not.
\usepackage{iftex}
\ifPDFTeX
  \usepackage[T1]{fontenc}
  \usepackage{lmodern}
\fi
\usepackage{microtype}
\usepackage{amsmath, amssymb, amsthm}
\usepackage{graphicx}
\usepackage{booktabs}
\usepackage[margin=1in]{geometry}
\usepackage{enumitem}
\usepackage{xcolor}
\usepackage{tcolorbox}

% Numbered statements. Theorems, lemmas and corollaries share one counter.
\theoremstyle{plain}
\newtheorem{theorem}{Theorem}[section]
\newtheorem{lemma}[theorem]{Lemma}
\newtheorem{corollary}[theorem]{Corollary}
\theoremstyle{definition}
\newtheorem{definition}[theorem]{Definition}
\newtheorem{example}[theorem]{Example}
\newtheorem*{keypoint}{Key point}
\theoremstyle{remark}
\newtheorem*{remark}{Remark}

% Draw every key point in a coloured box.
\tcolorboxenvironment{keypoint}{colback=blue!5, colframe=blue!50!black, boxrule=0.6pt, arc=2pt, left=6pt, right=6pt}

\numberwithin{equation}{section}

\usepackage[colorlinks=true, linkcolor=blue!60!black, citecolor=green!40!black, urlcolor=blue!60!black]{hyperref}

\title{Course Code: Course Title \\[0.5em] \large Lecture 1: Topic of the Lecture}
\author{Lecturer's Name}
\date{\today}

\begin{document}

\maketitle

\begin{abstract}
In this lecture we introduce the main objects of the course, prove the
first result about them, and work through an example. Before the lecture,
review your notes on the prerequisites listed in Section~\ref{sec:prereq}.
\end{abstract}

\tableofcontents

\section{Prerequisites}
\label{sec:prereq}

These notes assume you are comfortable with:
\begin{itemize}
  \item the material of the previous course in the sequence;
  \item basic algebra and proof by induction;
  \item summation notation.
\end{itemize}

\section{Definitions}
\label{sec:defs}

\begin{definition}
\label{def:sequence}
A \emph{sequence} is a function $a \colon \mathbb{N} \to \mathbb{R}$. We
write $a_n$ for its value at $n$. In these notes
$\mathbb{N} = \{0, 1, 2, \dots\}$ includes $0$.
\end{definition}

\begin{example}
\label{ex:arith}
The sequence $a_n = 2n + 1$ lists the odd numbers $1, 3, 5, \dots$
\end{example}

\section{Main Result}
\label{sec:main}

\begin{theorem}
\label{thm:sum}
For every natural number $n \geq 1$,
\begin{equation}
  \sum_{k=1}^{n} k = \frac{n(n+1)}{2} .
  \label{eq:sum}
\end{equation}
\end{theorem}

\begin{proof}
We use induction on $n$. For $n = 1$ both sides equal $1$. Assume
Equation~\eqref{eq:sum} holds for $n$. Then
\begin{align}
  \sum_{k=1}^{n+1} k &= \frac{n(n+1)}{2} + (n + 1) \\
  &= \frac{(n+1)(n+2)}{2} ,
\end{align}
which is the statement for $n + 1$.
\end{proof}

\begin{corollary}
\label{cor:odd}
The sum of the first $n$ odd numbers of Example~\ref{ex:arith} is $n^2$.
\end{corollary}

\begin{keypoint}
Induction proves a statement for every $n$ by checking a first case and
showing that each case implies the next. Theorem~\ref{thm:sum} is the
standard first example.
\end{keypoint}

\begin{remark}
There is also a direct proof: write the sum forwards and backwards and
add the two lines term by term.
\end{remark}

\section{Exercises}
\label{sec:exercises}

\begin{enumerate}[label=\arabic*.]
  \item Prove Corollary~\ref{cor:odd} by induction.
  \item Find a formula for $\sum_{k=1}^{n} k^2$ and prove it.
  \item Prove Theorem~\ref{thm:sum} using the direct argument in the remark.
\end{enumerate}

\end{document}
